%----------------------------------------------------------------------------------------------------%
%%%%%%%%%%%                                 MUSIK-BEFEHLE                                     %%%%%%%%%%%
%----------------------------------------------------------------------------------------------------%
% Optionale Ergänzung zur universalen newcommands.tex – nur für Musik-Dokumente.
% Einbinden NACH newcommands.tex:   %----------------------------------------------------------------------------------------------------%
%%%%%%%%%%%                                 MUSIK-BEFEHLE                                     %%%%%%%%%%%
%----------------------------------------------------------------------------------------------------%
% Optionale Ergänzung zur universalen newcommands.tex – nur für Musik-Dokumente.
% Einbinden NACH newcommands.tex:   %----------------------------------------------------------------------------------------------------%
%%%%%%%%%%%                                 MUSIK-BEFEHLE                                     %%%%%%%%%%%
%----------------------------------------------------------------------------------------------------%
% Optionale Ergänzung zur universalen newcommands.tex – nur für Musik-Dokumente.
% Einbinden NACH newcommands.tex:   %----------------------------------------------------------------------------------------------------%
%%%%%%%%%%%                                 MUSIK-BEFEHLE                                     %%%%%%%%%%%
%----------------------------------------------------------------------------------------------------%
% Optionale Ergänzung zur universalen newcommands.tex – nur für Musik-Dokumente.
% Einbinden NACH newcommands.tex:   \input{musik-befehle.tex}
% Getrennt gehalten, weil Musik-Dokumente zusätzlich musixtex brauchen (definiert \cup um)
% und die Tonnamen (\cis, \fis, ...) sonst leicht mit Mathe-Befehlen kollidieren.
% Quelle: https://github.com/mq90/latex-skills-schule  (CC BY 4.0, Martin Quandt)

%----------------------------------------------------------------------------------------------------%
%%%%%%%%%%%                              Musiktheoretische Abkürzungen                           %%%%%%%%%%
%----------------------------------------------------------------------------------------------------%

% Intervallbezeichnungen
\newcommand{\prim}{\text{Prime}}
\newcommand{\sek}{\text{Sekunde}}
\newcommand{\terz}{\text{Terz}}
\newcommand{\quart}{\text{Quarte}}
\newcommand{\quint}{\text{Quinte}}
\newcommand{\sext}{\text{Sexte}}
\newcommand{\sept}{\text{Septime}}
\newcommand{\okt}{\text{Oktave}}

% Akkordbezeichnungen
\newcommand{\dur}{\text{Dur}}
\newcommand{\moll}{\text{Moll}}

% Tonbezeichnungen (international)
\newcommand{\cis}{c$\sharp$}
\newcommand{\des}{d$\flat$}
\newcommand{\disTon}{d$\sharp$} % nicht \dis – das ist in newcommands.tex der Displaystyle
\newcommand{\es}{e$\flat$}
\newcommand{\fis}{f$\sharp$}
\newcommand{\ges}{g$\flat$}
\newcommand{\gis}{g$\sharp$}
\newcommand{\as}{a$\flat$}
\newcommand{\ais}{a$\sharp$}
\newcommand{\bb}{b$\flat$}
\newcommand{\bis}{b$\sharp$}

%----------------------------------------------------------------------------------------------------%
%%%%%%%%%%%                                Musik-spezifische Umgebungen                         %%%%%%%%%%
%----------------------------------------------------------------------------------------------------%

% Umgebung für Akkordfolgen
\newenvironment{akkordfolge}{%
    \begin{center}
    \begin{tabular}{|c|c|c|c|}
    \hline
    \textbf{Takt} & \textbf{1} & \textbf{2} & \textbf{3} \\
    \hline
    \textbf{Akkord} &
}{%
    \\
    \hline
    \end{tabular}
    \end{center}
}

% Umgebung für Tonleitern
\newenvironment{tonleiter}[1]{%
    \begin{center}
    \textbf{#1-Tonleiter}\\
    \begin{tabular}{|c|c|c|c|c|c|c|c|}
    \hline
    \textbf{Stufe} & I & II & III & IV & V & VI & VII \\
    \hline
    \textbf{Ton} &
}{%
    \\
    \hline
    \end{tabular}
    \end{center}
}

%----------------------------------------------------------------------------------------------------%
%%%%%%%%%%%                               MusiXTeX-spezifische Kürzel                           %%%%%%%%%%
%----------------------------------------------------------------------------------------------------%

% Notenwerte als Text
\newcommand{\ganze}{\textit{ganze Note}}
\newcommand{\halbe}{\textit{halbe Note}}
\newcommand{\viertel}{\textit{Viertelnote}}
\newcommand{\achtel}{\textit{Achtelnote}}
\newcommand{\sechzehntel}{\textit{Sechzehntelnote}}

% Pausenwerte als Text
\newcommand{\ganzpause}{\textit{ganze Pause}}
\newcommand{\halbpause}{\textit{halbe Pause}}
\newcommand{\viertelpause}{\textit{Viertelpause}}
\newcommand{\achtelpause}{\textit{Achtelpause}}
\newcommand{\sechzehntelpause}{\textit{Sechzehntelpause}}

% Häufige Taktarten
\newcommand{\viervierteltakt}{\ensuremath{\frac{4}{4}}\text{-Takt}}
\newcommand{\dreivierteltakt}{\ensuremath{\frac{3}{4}}\text{-Takt}}
\newcommand{\zweivierteltakt}{\ensuremath{\frac{2}{4}}\text{-Takt}}
\newcommand{\sechsachteltakt}{\ensuremath{\frac{6}{8}}\text{-Takt}}

%----------------------------------------------------------------------------------------------------%
%%%%%%%%%%%                                   Colorboxes für Musik                                %%%%%%%%%%
%----------------------------------------------------------------------------------------------------%

% Colorbox für Musiktheorie-Merksätze
\newtcolorbox{musiktheorie}{%
    lower separated=false,%
    colback=green!5!white,%
    colframe=green!75!black,%
    fonttitle=\bfseries,%
    title=Musiktheorie:,%
    enhanced,%
    breakable,%
    attach boxed title to top left={xshift=0.5cm,yshift=-2mm}%
}

% Colorbox für Übungen
\newtcolorbox{uebung}{%
    lower separated=false,%
    colback=orange!5!white,%
    colframe=orange!75!black,%
    fonttitle=\bfseries,%
    title=Übung:,%
    enhanced,%
    breakable,%
    attach boxed title to top left={xshift=0.5cm,yshift=-2mm}%
}

% Colorbox für Repertoire-Hinweise
\newtcolorbox{repertoire}{%
    lower separated=false,%
    colback=purple!5!white,%
    colframe=purple!75!black,%
    fonttitle=\bfseries,%
    title=Repertoire-Tipp:,%
    enhanced,%
    breakable,%
    attach boxed title to top left={xshift=0.5cm,yshift=-2mm}%
}

%----------------------------------------------------------------------------------------------------%
%%%%%%%%%%%                              Spezielle Musik-Befehle                                %%%%%%%%%%
%----------------------------------------------------------------------------------------------------%

% Befehl für Liedtitel mit Komponist
\newcommand{\liedtitel}[2]{%
    \begin{center}
    \Large\textbf{#1}\\
    \normalsize\textit{#2}
    \end{center}
    \vspace{0.5cm}
}

% Befehl für Tempo-Angaben
\newcommand{\tempo}[1]{%
    \textit{#1}
}

% Befehl für Tonart-Angaben
\newcommand{\tonart}[1]{%
    \textbf{Tonart:} #1
}

% Befehl für Taktart-Angaben
\newcommand{\taktart}[1]{%
    \textbf{Taktart:} #1
}

%----------------------------------------------------------------------------------------------------%
%%%%%%%%%%%                              Aufgaben-Umgebungen für Musik                          %%%%%%%%%%
%----------------------------------------------------------------------------------------------------%

% Umgebung für Gehörbildungs-Aufgaben
\newenvironment{gehoerbildung}[1]{%
    \begin{tcolorbox}[colback=blue!5!white,colframe=blue!50!black,title=Gehörbildung: #1]
}{%
    \end{tcolorbox}
}

% Umgebung für Rhythmus-Übungen
\newenvironment{rhythmus}[1]{%
    \begin{tcolorbox}[colback=red!5!white,colframe=red!50!black,title=Rhythmus-Übung: #1]
}{%
    \end{tcolorbox}
}

% Umgebung für Harmonielehre-Aufgaben
\newenvironment{harmonielehre}[1]{%
    \begin{tcolorbox}[colback=yellow!5!white,colframe=yellow!50!black,title=Harmonielehre: #1]
}{%
    \end{tcolorbox}
}

%----------------------------------------------------------------------------------------------------%
% Getrennt gehalten, weil Musik-Dokumente zusätzlich musixtex brauchen (definiert \cup um)
% und die Tonnamen (\cis, \fis, ...) sonst leicht mit Mathe-Befehlen kollidieren.
% Quelle: https://github.com/mq90/latex-skills-schule  (CC BY 4.0, Martin Quandt)

%----------------------------------------------------------------------------------------------------%
%%%%%%%%%%%                              Musiktheoretische Abkürzungen                           %%%%%%%%%%
%----------------------------------------------------------------------------------------------------%

% Intervallbezeichnungen
\newcommand{\prim}{\text{Prime}}
\newcommand{\sek}{\text{Sekunde}}
\newcommand{\terz}{\text{Terz}}
\newcommand{\quart}{\text{Quarte}}
\newcommand{\quint}{\text{Quinte}}
\newcommand{\sext}{\text{Sexte}}
\newcommand{\sept}{\text{Septime}}
\newcommand{\okt}{\text{Oktave}}

% Akkordbezeichnungen
\newcommand{\dur}{\text{Dur}}
\newcommand{\moll}{\text{Moll}}

% Tonbezeichnungen (international)
\newcommand{\cis}{c$\sharp$}
\newcommand{\des}{d$\flat$}
\newcommand{\disTon}{d$\sharp$} % nicht \dis – das ist in newcommands.tex der Displaystyle
\newcommand{\es}{e$\flat$}
\newcommand{\fis}{f$\sharp$}
\newcommand{\ges}{g$\flat$}
\newcommand{\gis}{g$\sharp$}
\newcommand{\as}{a$\flat$}
\newcommand{\ais}{a$\sharp$}
\newcommand{\bb}{b$\flat$}
\newcommand{\bis}{b$\sharp$}

%----------------------------------------------------------------------------------------------------%
%%%%%%%%%%%                                Musik-spezifische Umgebungen                         %%%%%%%%%%
%----------------------------------------------------------------------------------------------------%

% Umgebung für Akkordfolgen
\newenvironment{akkordfolge}{%
    \begin{center}
    \begin{tabular}{|c|c|c|c|}
    \hline
    \textbf{Takt} & \textbf{1} & \textbf{2} & \textbf{3} \\
    \hline
    \textbf{Akkord} &
}{%
    \\
    \hline
    \end{tabular}
    \end{center}
}

% Umgebung für Tonleitern
\newenvironment{tonleiter}[1]{%
    \begin{center}
    \textbf{#1-Tonleiter}\\
    \begin{tabular}{|c|c|c|c|c|c|c|c|}
    \hline
    \textbf{Stufe} & I & II & III & IV & V & VI & VII \\
    \hline
    \textbf{Ton} &
}{%
    \\
    \hline
    \end{tabular}
    \end{center}
}

%----------------------------------------------------------------------------------------------------%
%%%%%%%%%%%                               MusiXTeX-spezifische Kürzel                           %%%%%%%%%%
%----------------------------------------------------------------------------------------------------%

% Notenwerte als Text
\newcommand{\ganze}{\textit{ganze Note}}
\newcommand{\halbe}{\textit{halbe Note}}
\newcommand{\viertel}{\textit{Viertelnote}}
\newcommand{\achtel}{\textit{Achtelnote}}
\newcommand{\sechzehntel}{\textit{Sechzehntelnote}}

% Pausenwerte als Text
\newcommand{\ganzpause}{\textit{ganze Pause}}
\newcommand{\halbpause}{\textit{halbe Pause}}
\newcommand{\viertelpause}{\textit{Viertelpause}}
\newcommand{\achtelpause}{\textit{Achtelpause}}
\newcommand{\sechzehntelpause}{\textit{Sechzehntelpause}}

% Häufige Taktarten
\newcommand{\viervierteltakt}{\ensuremath{\frac{4}{4}}\text{-Takt}}
\newcommand{\dreivierteltakt}{\ensuremath{\frac{3}{4}}\text{-Takt}}
\newcommand{\zweivierteltakt}{\ensuremath{\frac{2}{4}}\text{-Takt}}
\newcommand{\sechsachteltakt}{\ensuremath{\frac{6}{8}}\text{-Takt}}

%----------------------------------------------------------------------------------------------------%
%%%%%%%%%%%                                   Colorboxes für Musik                                %%%%%%%%%%
%----------------------------------------------------------------------------------------------------%

% Colorbox für Musiktheorie-Merksätze
\newtcolorbox{musiktheorie}{%
    lower separated=false,%
    colback=green!5!white,%
    colframe=green!75!black,%
    fonttitle=\bfseries,%
    title=Musiktheorie:,%
    enhanced,%
    breakable,%
    attach boxed title to top left={xshift=0.5cm,yshift=-2mm}%
}

% Colorbox für Übungen
\newtcolorbox{uebung}{%
    lower separated=false,%
    colback=orange!5!white,%
    colframe=orange!75!black,%
    fonttitle=\bfseries,%
    title=Übung:,%
    enhanced,%
    breakable,%
    attach boxed title to top left={xshift=0.5cm,yshift=-2mm}%
}

% Colorbox für Repertoire-Hinweise
\newtcolorbox{repertoire}{%
    lower separated=false,%
    colback=purple!5!white,%
    colframe=purple!75!black,%
    fonttitle=\bfseries,%
    title=Repertoire-Tipp:,%
    enhanced,%
    breakable,%
    attach boxed title to top left={xshift=0.5cm,yshift=-2mm}%
}

%----------------------------------------------------------------------------------------------------%
%%%%%%%%%%%                              Spezielle Musik-Befehle                                %%%%%%%%%%
%----------------------------------------------------------------------------------------------------%

% Befehl für Liedtitel mit Komponist
\newcommand{\liedtitel}[2]{%
    \begin{center}
    \Large\textbf{#1}\\
    \normalsize\textit{#2}
    \end{center}
    \vspace{0.5cm}
}

% Befehl für Tempo-Angaben
\newcommand{\tempo}[1]{%
    \textit{#1}
}

% Befehl für Tonart-Angaben
\newcommand{\tonart}[1]{%
    \textbf{Tonart:} #1
}

% Befehl für Taktart-Angaben
\newcommand{\taktart}[1]{%
    \textbf{Taktart:} #1
}

%----------------------------------------------------------------------------------------------------%
%%%%%%%%%%%                              Aufgaben-Umgebungen für Musik                          %%%%%%%%%%
%----------------------------------------------------------------------------------------------------%

% Umgebung für Gehörbildungs-Aufgaben
\newenvironment{gehoerbildung}[1]{%
    \begin{tcolorbox}[colback=blue!5!white,colframe=blue!50!black,title=Gehörbildung: #1]
}{%
    \end{tcolorbox}
}

% Umgebung für Rhythmus-Übungen
\newenvironment{rhythmus}[1]{%
    \begin{tcolorbox}[colback=red!5!white,colframe=red!50!black,title=Rhythmus-Übung: #1]
}{%
    \end{tcolorbox}
}

% Umgebung für Harmonielehre-Aufgaben
\newenvironment{harmonielehre}[1]{%
    \begin{tcolorbox}[colback=yellow!5!white,colframe=yellow!50!black,title=Harmonielehre: #1]
}{%
    \end{tcolorbox}
}

%----------------------------------------------------------------------------------------------------%
% Getrennt gehalten, weil Musik-Dokumente zusätzlich musixtex brauchen (definiert \cup um)
% und die Tonnamen (\cis, \fis, ...) sonst leicht mit Mathe-Befehlen kollidieren.
% Quelle: https://github.com/mq90/latex-skills-schule  (CC BY 4.0, Martin Quandt)

%----------------------------------------------------------------------------------------------------%
%%%%%%%%%%%                              Musiktheoretische Abkürzungen                           %%%%%%%%%%
%----------------------------------------------------------------------------------------------------%

% Intervallbezeichnungen
\newcommand{\prim}{\text{Prime}}
\newcommand{\sek}{\text{Sekunde}}
\newcommand{\terz}{\text{Terz}}
\newcommand{\quart}{\text{Quarte}}
\newcommand{\quint}{\text{Quinte}}
\newcommand{\sext}{\text{Sexte}}
\newcommand{\sept}{\text{Septime}}
\newcommand{\okt}{\text{Oktave}}

% Akkordbezeichnungen
\newcommand{\dur}{\text{Dur}}
\newcommand{\moll}{\text{Moll}}

% Tonbezeichnungen (international)
\newcommand{\cis}{c$\sharp$}
\newcommand{\des}{d$\flat$}
\newcommand{\disTon}{d$\sharp$} % nicht \dis – das ist in newcommands.tex der Displaystyle
\newcommand{\es}{e$\flat$}
\newcommand{\fis}{f$\sharp$}
\newcommand{\ges}{g$\flat$}
\newcommand{\gis}{g$\sharp$}
\newcommand{\as}{a$\flat$}
\newcommand{\ais}{a$\sharp$}
\newcommand{\bb}{b$\flat$}
\newcommand{\bis}{b$\sharp$}

%----------------------------------------------------------------------------------------------------%
%%%%%%%%%%%                                Musik-spezifische Umgebungen                         %%%%%%%%%%
%----------------------------------------------------------------------------------------------------%

% Umgebung für Akkordfolgen
\newenvironment{akkordfolge}{%
    \begin{center}
    \begin{tabular}{|c|c|c|c|}
    \hline
    \textbf{Takt} & \textbf{1} & \textbf{2} & \textbf{3} \\
    \hline
    \textbf{Akkord} &
}{%
    \\
    \hline
    \end{tabular}
    \end{center}
}

% Umgebung für Tonleitern
\newenvironment{tonleiter}[1]{%
    \begin{center}
    \textbf{#1-Tonleiter}\\
    \begin{tabular}{|c|c|c|c|c|c|c|c|}
    \hline
    \textbf{Stufe} & I & II & III & IV & V & VI & VII \\
    \hline
    \textbf{Ton} &
}{%
    \\
    \hline
    \end{tabular}
    \end{center}
}

%----------------------------------------------------------------------------------------------------%
%%%%%%%%%%%                               MusiXTeX-spezifische Kürzel                           %%%%%%%%%%
%----------------------------------------------------------------------------------------------------%

% Notenwerte als Text
\newcommand{\ganze}{\textit{ganze Note}}
\newcommand{\halbe}{\textit{halbe Note}}
\newcommand{\viertel}{\textit{Viertelnote}}
\newcommand{\achtel}{\textit{Achtelnote}}
\newcommand{\sechzehntel}{\textit{Sechzehntelnote}}

% Pausenwerte als Text
\newcommand{\ganzpause}{\textit{ganze Pause}}
\newcommand{\halbpause}{\textit{halbe Pause}}
\newcommand{\viertelpause}{\textit{Viertelpause}}
\newcommand{\achtelpause}{\textit{Achtelpause}}
\newcommand{\sechzehntelpause}{\textit{Sechzehntelpause}}

% Häufige Taktarten
\newcommand{\viervierteltakt}{\ensuremath{\frac{4}{4}}\text{-Takt}}
\newcommand{\dreivierteltakt}{\ensuremath{\frac{3}{4}}\text{-Takt}}
\newcommand{\zweivierteltakt}{\ensuremath{\frac{2}{4}}\text{-Takt}}
\newcommand{\sechsachteltakt}{\ensuremath{\frac{6}{8}}\text{-Takt}}

%----------------------------------------------------------------------------------------------------%
%%%%%%%%%%%                                   Colorboxes für Musik                                %%%%%%%%%%
%----------------------------------------------------------------------------------------------------%

% Colorbox für Musiktheorie-Merksätze
\newtcolorbox{musiktheorie}{%
    lower separated=false,%
    colback=green!5!white,%
    colframe=green!75!black,%
    fonttitle=\bfseries,%
    title=Musiktheorie:,%
    enhanced,%
    breakable,%
    attach boxed title to top left={xshift=0.5cm,yshift=-2mm}%
}

% Colorbox für Übungen
\newtcolorbox{uebung}{%
    lower separated=false,%
    colback=orange!5!white,%
    colframe=orange!75!black,%
    fonttitle=\bfseries,%
    title=Übung:,%
    enhanced,%
    breakable,%
    attach boxed title to top left={xshift=0.5cm,yshift=-2mm}%
}

% Colorbox für Repertoire-Hinweise
\newtcolorbox{repertoire}{%
    lower separated=false,%
    colback=purple!5!white,%
    colframe=purple!75!black,%
    fonttitle=\bfseries,%
    title=Repertoire-Tipp:,%
    enhanced,%
    breakable,%
    attach boxed title to top left={xshift=0.5cm,yshift=-2mm}%
}

%----------------------------------------------------------------------------------------------------%
%%%%%%%%%%%                              Spezielle Musik-Befehle                                %%%%%%%%%%
%----------------------------------------------------------------------------------------------------%

% Befehl für Liedtitel mit Komponist
\newcommand{\liedtitel}[2]{%
    \begin{center}
    \Large\textbf{#1}\\
    \normalsize\textit{#2}
    \end{center}
    \vspace{0.5cm}
}

% Befehl für Tempo-Angaben
\newcommand{\tempo}[1]{%
    \textit{#1}
}

% Befehl für Tonart-Angaben
\newcommand{\tonart}[1]{%
    \textbf{Tonart:} #1
}

% Befehl für Taktart-Angaben
\newcommand{\taktart}[1]{%
    \textbf{Taktart:} #1
}

%----------------------------------------------------------------------------------------------------%
%%%%%%%%%%%                              Aufgaben-Umgebungen für Musik                          %%%%%%%%%%
%----------------------------------------------------------------------------------------------------%

% Umgebung für Gehörbildungs-Aufgaben
\newenvironment{gehoerbildung}[1]{%
    \begin{tcolorbox}[colback=blue!5!white,colframe=blue!50!black,title=Gehörbildung: #1]
}{%
    \end{tcolorbox}
}

% Umgebung für Rhythmus-Übungen
\newenvironment{rhythmus}[1]{%
    \begin{tcolorbox}[colback=red!5!white,colframe=red!50!black,title=Rhythmus-Übung: #1]
}{%
    \end{tcolorbox}
}

% Umgebung für Harmonielehre-Aufgaben
\newenvironment{harmonielehre}[1]{%
    \begin{tcolorbox}[colback=yellow!5!white,colframe=yellow!50!black,title=Harmonielehre: #1]
}{%
    \end{tcolorbox}
}

%----------------------------------------------------------------------------------------------------%
% Getrennt gehalten, weil Musik-Dokumente zusätzlich musixtex brauchen (definiert \cup um)
% und die Tonnamen (\cis, \fis, ...) sonst leicht mit Mathe-Befehlen kollidieren.
% Quelle: https://github.com/mq90/latex-skills-schule  (CC BY 4.0, Martin Quandt)

%----------------------------------------------------------------------------------------------------%
%%%%%%%%%%%                              Musiktheoretische Abkürzungen                           %%%%%%%%%%
%----------------------------------------------------------------------------------------------------%

% Intervallbezeichnungen
\newcommand{\prim}{\text{Prime}}
\newcommand{\sek}{\text{Sekunde}}
\newcommand{\terz}{\text{Terz}}
\newcommand{\quart}{\text{Quarte}}
\newcommand{\quint}{\text{Quinte}}
\newcommand{\sext}{\text{Sexte}}
\newcommand{\sept}{\text{Septime}}
\newcommand{\okt}{\text{Oktave}}

% Akkordbezeichnungen
\newcommand{\dur}{\text{Dur}}
\newcommand{\moll}{\text{Moll}}

% Tonbezeichnungen (international)
\newcommand{\cis}{c$\sharp$}
\newcommand{\des}{d$\flat$}
\newcommand{\disTon}{d$\sharp$} % nicht \dis – das ist in newcommands.tex der Displaystyle
\newcommand{\es}{e$\flat$}
\newcommand{\fis}{f$\sharp$}
\newcommand{\ges}{g$\flat$}
\newcommand{\gis}{g$\sharp$}
\newcommand{\as}{a$\flat$}
\newcommand{\ais}{a$\sharp$}
\newcommand{\bb}{b$\flat$}
\newcommand{\bis}{b$\sharp$}

%----------------------------------------------------------------------------------------------------%
%%%%%%%%%%%                                Musik-spezifische Umgebungen                         %%%%%%%%%%
%----------------------------------------------------------------------------------------------------%

% Umgebung für Akkordfolgen
\newenvironment{akkordfolge}{%
    \begin{center}
    \begin{tabular}{|c|c|c|c|}
    \hline
    \textbf{Takt} & \textbf{1} & \textbf{2} & \textbf{3} \\
    \hline
    \textbf{Akkord} &
}{%
    \\
    \hline
    \end{tabular}
    \end{center}
}

% Umgebung für Tonleitern
\newenvironment{tonleiter}[1]{%
    \begin{center}
    \textbf{#1-Tonleiter}\\
    \begin{tabular}{|c|c|c|c|c|c|c|c|}
    \hline
    \textbf{Stufe} & I & II & III & IV & V & VI & VII \\
    \hline
    \textbf{Ton} &
}{%
    \\
    \hline
    \end{tabular}
    \end{center}
}

%----------------------------------------------------------------------------------------------------%
%%%%%%%%%%%                               MusiXTeX-spezifische Kürzel                           %%%%%%%%%%
%----------------------------------------------------------------------------------------------------%

% Notenwerte als Text
\newcommand{\ganze}{\textit{ganze Note}}
\newcommand{\halbe}{\textit{halbe Note}}
\newcommand{\viertel}{\textit{Viertelnote}}
\newcommand{\achtel}{\textit{Achtelnote}}
\newcommand{\sechzehntel}{\textit{Sechzehntelnote}}

% Pausenwerte als Text
\newcommand{\ganzpause}{\textit{ganze Pause}}
\newcommand{\halbpause}{\textit{halbe Pause}}
\newcommand{\viertelpause}{\textit{Viertelpause}}
\newcommand{\achtelpause}{\textit{Achtelpause}}
\newcommand{\sechzehntelpause}{\textit{Sechzehntelpause}}

% Häufige Taktarten
\newcommand{\viervierteltakt}{\ensuremath{\frac{4}{4}}\text{-Takt}}
\newcommand{\dreivierteltakt}{\ensuremath{\frac{3}{4}}\text{-Takt}}
\newcommand{\zweivierteltakt}{\ensuremath{\frac{2}{4}}\text{-Takt}}
\newcommand{\sechsachteltakt}{\ensuremath{\frac{6}{8}}\text{-Takt}}

%----------------------------------------------------------------------------------------------------%
%%%%%%%%%%%                                   Colorboxes für Musik                                %%%%%%%%%%
%----------------------------------------------------------------------------------------------------%

% Colorbox für Musiktheorie-Merksätze
\newtcolorbox{musiktheorie}{%
    lower separated=false,%
    colback=green!5!white,%
    colframe=green!75!black,%
    fonttitle=\bfseries,%
    title=Musiktheorie:,%
    enhanced,%
    breakable,%
    attach boxed title to top left={xshift=0.5cm,yshift=-2mm}%
}

% Colorbox für Übungen
\newtcolorbox{uebung}{%
    lower separated=false,%
    colback=orange!5!white,%
    colframe=orange!75!black,%
    fonttitle=\bfseries,%
    title=Übung:,%
    enhanced,%
    breakable,%
    attach boxed title to top left={xshift=0.5cm,yshift=-2mm}%
}

% Colorbox für Repertoire-Hinweise
\newtcolorbox{repertoire}{%
    lower separated=false,%
    colback=purple!5!white,%
    colframe=purple!75!black,%
    fonttitle=\bfseries,%
    title=Repertoire-Tipp:,%
    enhanced,%
    breakable,%
    attach boxed title to top left={xshift=0.5cm,yshift=-2mm}%
}

%----------------------------------------------------------------------------------------------------%
%%%%%%%%%%%                              Spezielle Musik-Befehle                                %%%%%%%%%%
%----------------------------------------------------------------------------------------------------%

% Befehl für Liedtitel mit Komponist
\newcommand{\liedtitel}[2]{%
    \begin{center}
    \Large\textbf{#1}\\
    \normalsize\textit{#2}
    \end{center}
    \vspace{0.5cm}
}

% Befehl für Tempo-Angaben
\newcommand{\tempo}[1]{%
    \textit{#1}
}

% Befehl für Tonart-Angaben
\newcommand{\tonart}[1]{%
    \textbf{Tonart:} #1
}

% Befehl für Taktart-Angaben
\newcommand{\taktart}[1]{%
    \textbf{Taktart:} #1
}

%----------------------------------------------------------------------------------------------------%
%%%%%%%%%%%                              Aufgaben-Umgebungen für Musik                          %%%%%%%%%%
%----------------------------------------------------------------------------------------------------%

% Umgebung für Gehörbildungs-Aufgaben
\newenvironment{gehoerbildung}[1]{%
    \begin{tcolorbox}[colback=blue!5!white,colframe=blue!50!black,title=Gehörbildung: #1]
}{%
    \end{tcolorbox}
}

% Umgebung für Rhythmus-Übungen
\newenvironment{rhythmus}[1]{%
    \begin{tcolorbox}[colback=red!5!white,colframe=red!50!black,title=Rhythmus-Übung: #1]
}{%
    \end{tcolorbox}
}

% Umgebung für Harmonielehre-Aufgaben
\newenvironment{harmonielehre}[1]{%
    \begin{tcolorbox}[colback=yellow!5!white,colframe=yellow!50!black,title=Harmonielehre: #1]
}{%
    \end{tcolorbox}
}

%----------------------------------------------------------------------------------------------------%